% !TeX root = ../main.tex

\begin{denotation}[3cm]

\item[IENF]{
  表皮內神經纖維 (Intraepidermal Nerve Fiber)
}

\item[IENFD]{
  表皮內神經纖維密度 (Intraepidermal Nerve Fiber Density)
}

\item[SFN]{
  小纖維神經病變 (Small Fiber Neuropathy)
}

\item[DEJ]{
  表皮／真皮交界 (Dermal--Epidermal Junction)
}

\item[AEL]{
  標註擴張連結 (Annotation Expansion Linker)，本論文所提出之方法
}

\item[ROI]{
  感興趣區域 (Region of Interest)
}

\item[CLAHE]{
  對比受限直方圖均衡化 (Contrast Limited Adaptive Histogram Equalization)
}

\item[MST]{
  最小生成樹 (Minimum Spanning Tree)
}

\item[MSF]{
  最小生成森林 (Minimum Spanning Forest)
}

\item[CC]{
  連通元件 (Connected Component)
}

\item[$I$]{
  輸入之組織切片影像 (Input histological image)
}

\item[$I_{g}$]{
  綠色通道影像 (Green channel image)
}

\item[$M$]{
  表皮遮罩 (Epidermis mask)
}

\item[$M_{\text{ROI}}$]{
  向下膨脹後之感興趣區域遮罩 (Vertically dilated ROI mask)
}

\item[$\text{offset}$]{
  感興趣區域向下膨脹之圓形結構元半徑（預設 $50$ 像素）
}

\item[$A$]{
  二值化之纖維微粒標註影像 (Binary particle annotation)
}

\item[$I_{\text{th}}$]{
  白色頂帽去背景後之影像 (White top-hat background-removed image)
}

\item[$r_{\text{bg}}$]{
  白色頂帽結構元半徑（預設 $2$ 像素）
}

\item[$I_{\text{eq}}$]{
  CLAHE 局部對比增強後之影像 (CLAHE-equalized image)
}

\item[$\beta$]{
  CLAHE 之對比裁剪閾值（預設 $30$）
}

\item[$n_{\text{tile}}$]{
  CLAHE 之方塊尺度（預設 $768$ 像素）
}

\item[$\sigma_{\min},\,\sigma_{\max}$]{
  Sato 多尺度線狀結構濾波之最小／最大尺度（取 $1$／$4$ 像素）
}

\item[$I_{\text{enhanced}}$]{
  Sato 線狀結構濾波並正規化後之纖維響應影像 (Vesselness-enhanced image)
}

\item[$C$]{
  像素成本圖 (Pixel-level cost map)，$C(p) = \exp\!\bigl(1 - \tilde{I}_{\text{enhanced}}(p)\bigr) - 1$
}

\item[$L$]{
  連通元件標記圖 (Connected component label map)
}

\item[$N$]{
  標註中連通元件之總數 (Number of annotation components)
}

\item[$D(p)$]{
  像素 $p$ 從任一種子經最小成本路徑所累積之成本 (Accumulated cost)
}

\item[$O(p)$]{
  像素 $p$ 所歸屬之元件索引 (Owner component of pixel $p$)
}

\item[$\mathcal{M}_{ab}$]{
  元件對 $(a, b)$ 之候選相遇點集合 (Candidate meeting points)
}

\item[$w_{ab}$]{
  元件對 $(a, b)$ 之最小相遇成本，即元件層圖之邊權 (Edge weight)
}

\item[$G_{c}$]{
  元件層圖 (Component-level graph)，節點為連通元件、邊為跨元件相遇點
}

\item[$\tau$]{
  相鄰元件間相遇成本之裁剪閾值（預設 $20$）
}

\item[$G_{\text{mst}}$]{
  由 $G_{c}$ 經裁剪後求得之最小生成森林，每棵樹對應一條重建纖維
}

\item[$\Pi_{(a,b)}$]{
  最小生成森林邊 $(a, b)$ 經前驅回溯所得之像素路徑 (Back-tracked pixel path)
}

\item[$\mathcal{B}$]{
  彙整所有回溯路徑而成之橋接遮罩 (Bridge mask)
}

\item[$r_{d}$]{
  骨架化前橋接路徑之膨脹半徑（預設 $1$ 像素）
}

\item[$G_{\text{skel}}$]{
  橋接遮罩與原標註聯集後骨架化所得之骨架圖 (Skeleton graph)
}

\item[$\ell_{\min}$]{
  短樁修剪之最小線段長度；長度小於此值之懸掛短樁將被移除（預設 $3$ 像素）
}

\item[$G_{\text{trim}}$]{
  經分段與短樁修剪後之最終像素級拓樸圖 (Trimmed topology graph)
}

\item[$N_{\min}$]{
  一條獨立神經須涵蓋之最小纖維微粒數（預設 $1$）
}

\item[$G_{\text{cov}}$]{
  經微粒涵蓋篩選後保留之拓樸圖 (Coverage-filtered graph)
}

\item[$N_{\text{cross}}$]{
  AEL 所輸出之有效跨越數 (Effective crossing count)
}

\item[$N^{*}_{\text{cross}}$]{
  專家依 EFNS 規則人工計數所得之有效跨越數真值
}

\item[$S_{\text{pred}}$]{
  演算法輸出之重建纖維骨架 (Predicted skeleton)
}

\item[$S_{\text{gt}}$]{
  專家標註之纖維骨架真值 (Ground-truth skeleton)
}

\item[$r_{v}$]{
  拓樸品質評估時骨架之容忍膨脹半徑（取 $1.28\,\mu\text{m}$）
}

\end{denotation}
